\subsection{RELATIONS WITH COXETER SYSTEMS}

THEOREM 2. The pair \((\mathrm{W},\mathrm{S})\) is a Coxeter system. Moreover, for \(s\in \mathrm{S}\) and \(w\in \mathrm{W}\) , the relations \(\mathrm{C}(sw) = \mathrm{C}(s).C(w)\) and \(l_{\mathbb{S}}(sw) > l_{\mathbb{S}}(w)\) are equivalent.

For any \(s \in S\) , let \(P_s\) be the set of elements \(w \in W\) such that

\[
\mathrm{C} (s). \mathrm{C} (w) = \mathrm{C} (s w).
\]

We are going to verify that the \(P_{s}\) satisfy conditions ( \(A'\) ), ( \(B'\) ) and (C) of §1, no. 7; the two assertions of the theorem will then follow from Prop. 6 of §1, no. 7.

Condition \((\mathbf{A}^{\prime})\) is clear.

We verify (B'). If \(P_{s}\) and \(sP_{s}\) had an element w in common, we would have \(w \in P_{s}\) and \(sw \in P_{s}\) , and hence

\[
\mathrm{C} (s). \mathrm{C} (w) = \mathrm{C} (s w), \quad \mathrm{C} (s). \mathrm{C} (s w) = \mathrm{C} (w).
\]

It would follow that \(\mathrm{C}(s).\mathrm{C}(s).\mathrm{C}(w) = \mathrm{C}(w)\) and, by formula (5), this would imply that \(\mathrm{C}(w) = \mathrm{C}(sw)\) , which would contradict Th. 1.

We verify (C). Let \(s, s' \in S\) and \(w, w' \in W\) with \(w' = ws'\) . The assumption that \(w \in P_s\) and \(w' \notin P_s\) implies that

\[
\mathrm{C} (s w) = \mathrm{C} (s). \mathrm{C} (w)\tag{9}
\]

\[
\mathrm{C} (w ^ {\prime}) \subset \mathrm{C} (s). \mathrm{C} (w ^ {\prime})\tag{10}
\]

by (3).

From (9) and the relation \(w = w's'\) , it follows that

\[
\mathrm{C} (s) w ^ {\prime} s ^ {\prime} \mathrm{B} = \mathrm{C} (s w).\tag{11}
\]

By formula \((2')\) , \(C(w')\) . \(C(s') \subset C(w') \cup C(w's')\) , which immediately implies that

\[
\mathrm{C} (w ^ {\prime}) s ^ {\prime} \mathrm{B} \subset \mathrm{C} (w s ^ {\prime}) \cup \mathrm{C} (w).\tag{12}
\]

Since \(C(w')\) is a union of left cosets gB and since

\[
\mathrm{C} (s). \mathrm{C} (w ^ {\prime}) = \mathrm{C} (s) w ^ {\prime} \mathrm{B},
\]

formula (10) shows that \(C(s)w'\) meets \(C(w')\) and a fortiori that \(C(s)w's'B\) meets \(C(w')s'B\) . It follows from formulas (11) and (12) that the double coset \(C(sw)\) is equal to one of the double cosets \(C(ws')\) and \(C(w)\) ; since \(sw \neq w\) , Th. 1 allows us to conclude that \(sw = ws'\) .

COROLLARY 1. Let \(w_{1}, \ldots, w_{q} \in \mathbf{W}\) and let \(w = w_{1} \ldots w_{q}\) . If

\[
l _ {S} (w) = l _ {S} \left(w _ {1}\right) + \dots + l _ {S} \left(w _ {q}\right),
\]

then

\[
\mathrm{C} (w) = \mathrm{C} (w _ {1}) \dots \mathrm{C} (w _ {q}).
\]

On taking reduced decompositions of the \(w_{i}\) , one is reduced to the case of a reduced decomposition

\[
w = s _ {1} \dots s _ {q}, \quad \text { with } s _ {i} \in S.
\]

If \(u = s_2 \ldots s_q\) , then \(w = s_1u\) and \(l_{\mathbb{S}}(s_1u) > l_{\mathbb{S}}(u)\) , so \(\mathrm{C}(w) = \mathrm{C}(s_1). \mathrm{C}(u)\) by the theorem. The required formula follows from this by induction on \(q\) .

COROLLARY 2. Let \(w \in W\) and let \(T_w\) be the subset of \(W\) associated to \(w\) by the procedure of Lemma 2 of § 1, no. 4. If \(t \in T_w\) , then

\[
\mathrm{C} (t) \subset \mathrm{C} (w). \mathrm{C} (w ^ {- 1}).
\]

If \(t \in \mathrm{T}_w\) , there exist by definition elements \(w', w'' \in \mathrm{W}\) and \(s \in \mathrm{S}\) such that

\[
w = w ^ {\prime} s w ^ {\prime \prime}, l _ {S} (w) = l _ {S} (w ^ {\prime}) + l _ {S} (w ^ {\prime \prime}) + 1 \quad \text {and} \quad t = w ^ {\prime} s w ^ {\prime - 1}.
\]

By Cor. 1,

\[
\mathrm{C} (w). \mathrm{C} (w ^ {- 1}) = \mathrm{C} (w ^ {\prime}). \mathrm{C} (s). \mathrm{C} (w ^ {\prime \prime}). \mathrm{C} (w ^ {\prime \prime - 1}). \mathrm{C} (s). \mathrm{C} (w ^ {\prime - 1}).
\]

Hence,

\[
\mathrm{C} (w). \mathrm{C} (w ^ {- 1}) \supset \mathrm{C} (w ^ {\prime}). \mathrm{C} (s). \mathrm{C} (s). \mathrm{C} (w ^ {\prime - 1}).
\]

By (4), \(\mathrm{C}(s) \subset \mathrm{C}(s). \mathrm{C}(s)\) . Hence,

\[
\mathrm{C} (w). \mathrm{C} (w ^ {- 1}) \supset \mathrm{C} (w ^ {\prime}). \mathrm{C} (s). \mathrm{C} (w ^ {\prime - 1}) \supset \mathrm{C} (t).
\]

COROLLARY 3. Let \(w \in W\) and let \(H_w\) be the subgroup of \(G\) generated by \(C(w).C(w^{-1})\) . Then:

\begin{enumerate}
\def\labelenumi{\alph{enumi})}
\tightlist
\item
  For any reduced decomposition \((s_1,\ldots ,s_q)\) of \(w\) ,
\end{enumerate}

\[
\mathrm{C} (s _ {j}) \subset \mathrm{H} _ {w} \quad \text { for } 1 \leqslant j \leqslant q.
\]

\begin{enumerate}
\def\labelenumi{\alph{enumi})}
\setcounter{enumi}{1}
\tightlist
\item
  The group \(\mathrm{H}_w\) contains \(\mathrm{C}(w)\) and is generated by \(\mathrm{C}(w)\) .
\end{enumerate}

We prove \(a)\) by induction on \(j\) . Assume that \(C(s_k)\) is contained in \(H_w\) for \(k < j\) . Let

\[
t = (s _ {1} \dots s _ {j - 1}) s _ {j} (s _ {1} \dots s _ {j - 1}) ^ {- 1}.
\]

The element \(t\) belongs to the subset \(\mathrm{T}_w\) of \(\mathbf{W}\) defined in Lemma 2 of §1, no. 4. By Cor. 2, \(\mathrm{C}(t) \subset \mathrm{H}_w\) , and hence \(\mathrm{C}(s_j) \subset \mathrm{H}_w\) .

Since \(\mathrm{C}(w) = \mathrm{C}(s_1) \ldots \mathrm{C}(s_q)\) , cf.~Cor. 1, we have \(\mathrm{C}(w) \subset \mathrm{H}_w\) , and \(b)\) follows.

Example. Th. 2, applied to the Tits system described in no. 2, shows that the symmetric group \(S_{n}\) , with the set of transpositions of consecutive elements, is a Coxeter group.

